% Dueño: B. A24, A25: una conclusión por cada resultado de recuperación y archivos sin codificar.
La evaluación de la autocorrelación del cepstro en el dominio espectral a partir del logaritmo de la magnitud espectral y la normalización por valor eficaz elimina la dependencia de escala de la señal portadora y las discontinuidades de fase, permitiendo distinguir de forma robusta los retardos $d_0$ y $d_1$ frente al piso cepstral de ventanas sin eco.

La tasa de recuperación crece monótonamente con la amplitud del eco y con el factor de tamaño de ventana $K$; la configuración por defecto ($t_0 = \SI{1000}{\micro\second}$, $t_1 = \SI{1500}{\micro\second}$, $a = \SI{30}{\percent}$ y $K = 25$) alcanza una recuperación media de hasta \SI{93.07}{\percent} a una tasa binaria de \SI{26.73}{bps}, cumpliendo el requerimiento mínimo del \SI{90}{\percent}.

El contenido espectral del género musical incide directamente en el desempeño del decodificador: las señales con espectro denso o percusivo (\textit{ragtime}, \textit{jazz} y electrónica) superan el \SI{91.9}{\percent} de recuperación con $a = \SI{30}{\percent}$, mientras que las cuerdas frotadas sostenidas presentan mayor interferencia cepstral cuando sus periodos armónicos coinciden con los retardos de codificación.

El criterio de umbral sobre $\max(|R_{cc}(d_0)|, |R_{cc}(d_1)|)$ combinado con la validación de actividad mínima por señal permite identificar entre el \SI{82.77}{\percent} y el \SI{94.01}{\percent} de las ventanas individuales sin eco y garantiza una tasa de falsos positivos nula (cero bits y cero metadatos espurios) al procesar archivos de audio no codificados.
